\chapter{Methodology}
\label{ch:methodology}

\section{Synthetic Supply Chain Data Model}
\label{sec:data_model}

The data model used in this project represents a fictional but structurally realistic
supply chain covering electronic components and consumer hardware. Ten event types drive
the simulation: sales (product sold from a warehouse in a region), restocks (units added
to a warehouse from a supplier), shipments (in-transit orders with an ETA), returns
(products sent back with a reason code), price changes (cost or sale price adjustments),
supplier delays (ETA extensions with a severity code), quality issues (defect rates above
a threshold, assigned a severity), demand spikes (relative increase in regional demand
for a product), stockouts (warehouse inventory reaching zero with pending orders), and
new contracts (supplier-product agreements with a total value).

The fixed topology establishes the entities. Twenty product nodes cover categories like
RAM modules, GPU cards, SSDs, laptops, and network hardware. Six supplier nodes
(ShenzhenDirect, TechSource Asia, BerlinElectronics, EuroComponents GmbH, AmeriParts Inc,
MumbaiMicro Ltd, TokyoTech Corp) vary in reliability and geographic coverage. Five
warehouse nodes (WH-LON, WH-NYC, WH-BER, WH-TKY, WH-SHZ) represent major logistics
hubs. Four region nodes (Europe, North America, Asia Pacific, Latin America) group the
warehouses by geography.

At steady state after processing a representative event batch, the knowledge graph
contains 30 nodes and 368 edges. The edge count is higher than a simple product count
would suggest because each product can be supplied by multiple suppliers, stored in
multiple warehouses, and sold into multiple regions. The density of the graph is what
makes centrality measures informative: a high-betweenness-centrality supplier is one
whose relationships span many products and warehouses, meaning its disruptions propagate
broadly.

\section{Knowledge Graph Construction}
\label{sec:graph_construction}

The Graph Builder processes event batches in two phases. In the update phase, each event
modifies the graph according to its type: a sale decrements the inventory quantity
attribute on the relevant warehouse-product edge, a quality issue sets the
\texttt{anomaly\_score} attribute on the affected product node to a severity-weighted
value, and a new contract adds an edge between the supplier node and the product node with
a \texttt{contract\_value} attribute. In the snapshot phase, the current graph state is
deep-copied before the batch begins and again after it ends, giving the delta engine two
comparable snapshots.

Node attributes are typed and bounded. The \texttt{anomaly\_score} attribute on product
and supplier nodes is a float in $[0, 1]$; it is set to 0.3 for minor quality issues,
0.6 for major issues, and 0.9 for critical issues. The \texttt{reliability} attribute on
supplier nodes is also a float in $[0, 1]$, decremented by 0.1 for each delay event and
incremented by 0.05 for each on-time delivery. Edge attributes include \texttt{quantity}
(current inventory), \texttt{pending\_orders}, \texttt{defect\_rate}, and
\texttt{contract\_value} where applicable.

The graph is implemented using NetworkX's \texttt{DiGraph} class, which supports directed
edges and arbitrary attribute dictionaries on both nodes and edges. The choice of a
directed graph reflects the directional nature of supply chain relationships: a ``supplies''
edge runs from a supplier node to a product node, not the reverse. Centrality measures are
computed on the undirected projection of the graph (treating all edges as undirected) for
the betweenness calculation, since the relevant question is whether a node sits on paths
between other nodes regardless of direction.

\section{Semantic Delta Computation}
\label{sec:semantic_delta}

The delta engine takes the before and after snapshots and produces a structured diff.
New nodes are any nodes present in the after-snapshot but absent in the before-snapshot,
indicating new suppliers, warehouses, or products added during the batch. Removed nodes
are the converse. Changed attribute nodes are nodes present in both snapshots whose
attribute dictionary differs by at least one key-value pair. New and removed edges follow
the same logic.

The delta representation is structured as a Python dictionary with five lists, one per
category. Each item in a list carries the node or edge identifier, the relevant attributes
(for changed-attribute items: the old value and new value of each changed attribute), and
a pre-computed reference to the node's position in the graph (its neighbors, its centrality
rank, its node type). Including this graph context in the delta item is important because
the saliency scorer needs it to compute centrality-based scores without re-querying the
full graph for every item.

A delta computed on a 1,000-event batch typically contains 15--40 distinct items, depending
on the volatility of the batch. Low-volatility batches with routine sales and restocks
produce few attribute changes above the detection threshold. High-volatility batches with
quality alerts, demand spikes, and supplier delays can produce 60--80 items before saliency
filtering. The delta computation itself is $O(n_{\text{nodes}} + n_{\text{edges}})$ in the
graph size, which at 30 nodes and 368 edges runs in well under a millisecond.

\section{Saliency Scoring}
\label{sec:saliency_scoring}

Each delta item receives a saliency score computed by the formula:

\begin{equation}
\text{saliency}(v) = \alpha \cdot d(v) + \beta \cdot b(v) + \gamma \cdot m(v)
\label{eq:saliency}
\end{equation}

where $d(v)$ is the degree centrality of the primary node involved in the change,
$b(v)$ is the betweenness centrality of that node (both normalized to $[0, 1]$ over
the current graph), and $m(v)$ is the normalized magnitude of the attribute change.
For attribute changes, $m(v)$ is the absolute value of the new-minus-old difference
divided by the maximum observed magnitude for that attribute type in the current batch.
For structural changes (new or removed nodes and edges), $m(v)$ is set to 0.5 as a
fixed structural importance weight.

The mixing coefficients $\alpha = 0.35$, $\beta = 0.35$, and $\gamma = 0.30$ were
calibrated on the ground truth evaluation set. A grid search over the three coefficients
at 0.05 increments found that this combination maximized recall on the four ground truth
patterns while keeping the fraction of events passing the threshold below 20\%, which
is the target compression ratio needed to achieve 5$\times$ or better compression at
the 100-event scale.

The threshold $\theta = 0.4$ is applied after scoring. Items with saliency $< 0.4$
are labeled NOISE and discarded. Items with $0.4 \leq \text{saliency} \leq 0.5$ are
INFORMATIONAL. Items with $0.5 < \text{saliency} \leq 0.75$ are IMPORTANT. Items with
$\text{saliency} > 0.75$ are CRITICAL. Table~\ref{tab:severity_levels} summarizes
the severity classification scheme.

\begin{table}[H]
\centering
\caption{Distillation severity levels and agent routing}
\label{tab:severity_levels}
\begin{tabular}{llll}
\toprule
\textbf{Level} & \textbf{Threshold} & \textbf{Included in Briefing} & \textbf{Agent Routing} \\
\midrule
CRITICAL      & $> 0.75$       & Yes & All three agents \\
IMPORTANT     & $0.5$--$0.75$  & Yes & Analyst + Risk \\
INFORMATIONAL & $0.4$--$0.5$   & Yes (summary only) & Analyst only \\
NOISE         & $< 0.4$        & No  & Discarded \\
\bottomrule
\end{tabular}
\end{table}

\section{Distillation Format}
\label{sec:distillation_format}

The distiller takes the scored delta items and assembles them into a Business Intelligence
Briefing. The briefing has a fixed structure: a header with metadata (batch timestamp,
event count, graph node and edge counts), a CRITICAL section listing the highest-saliency
items with full relationship context, an IMPORTANT section listing mid-saliency items
with condensed context, and an INFORMATIONAL section giving a one-line summary of each
low-saliency surviving item.

The relationship context included for each item is the key differentiator from aggregation
approaches. For a CRITICAL item flagging a quality issue at ShenzhenDirect affecting
GPU Card RTX, the briefing includes: the supplier's current reliability score, the list
of products it supplies, the warehouses where those products are stored, the pending
order counts at each warehouse, and the delta values for the anomaly score and defect
rate attributes. A pure aggregation approach that says ``ShenzhenDirect: 1 critical quality
issue'' provides the label but not the blast radius, which is exactly the information an
analyst agent needs to assess downstream impact.

The target briefing length is 600--900 tokens. In practice, briefings for 1,000-event
batches average 734 tokens (from the ablation results), within this range. The
briefing is plain text with no JSON nesting, because early experiments showed that heavily
structured JSON briefings cost more tokens per unit of information than prose with inline
attribute values.

\section{Evaluation Methodology}
\label{sec:evaluation}

The evaluation covers three separate measurements. The compression and cost evaluation
measures token counts for naive, aggregation, and pipeline methods at five scales
(100, 300, 500, 1,000, and 3,000 events) and derives cost savings by assuming a fixed
per-token price. The quality evaluation measures category recall by having a human
assessor judge whether each method's LLM response correctly identifies eight analytical
categories (supplier reliability, stockout detection, demand spike, quality issue,
price volatility, supply chain risk, warehouse anomaly, and actionable recommendations).
The ground truth evaluation injects four known anomaly patterns into a controlled batch
and checks whether each method's LLM response mentions the injected entities and
relationship structures.

The quality evaluation methodology is post-hoc analysis of saved LLM responses rather
than a live re-run for each measurement, which is a limitation acknowledged in
Chapter~\ref{ch:conclusion}. The ground truth evaluation runs live against the actual
pipeline, injecting the patterns into the data generator and verifying detection.
The ablation study is also live: the benchmark runner generates the same event batch
and feeds it to all three methods, recording token counts and pipeline timing for each.
